% !TEX root = ../monograph.tex
\chapter{The renormalization transformation and the effective density}
\label{ch04:blocking}

A renormalization transformation, called a step, is first an exact integration, and only then an
estimate. We define that integration on every fine configuration, including
large fields. Its output is generally not a Wilson density. We therefore
also introduce a representation of the effective density: a prescribed leading
action, a normalization constant, localized terms on small-field regions,
and a sum over polymers covering the large-field sites.

We use the gauge group and Haar measure of Chapter~\ref{ch01:lattice}
and the normalized matrix norm $\|\cdot\|_G$ of \eqref{ch01:matrix-norm}.
Left or right multiplication by a group element preserves this norm.
Recall that $a_m=D_m^{-1}$ is the lattice
spacing and $R_p$ the period, both in continuum units, and the odd ratio $L_m$ relates
the scales with indices $m+1$ and $m$. In this chapter we examine a single step
$m+1\to m$; unlike in~\cite{Balaban1984Propagators}, a step is not followed by a rescaling to the unit lattice. Write $L=L_m$, fix a positive integer \emph{smoothing order} $j$ (the number of convolution factors in the averaging weight $w^{\rm av}$ of Section~\ref{ch04:kernel-section}),
and set $R=(j+3)L$. Fine links $U$ live on $\Lambda_{R_pD_{m+1}}$ and coarse
links $V$ on $\Lambda_{R_pD_m}$. Patches are first described on the infinite
lattice $\ZZ^4$ and then placed periodically. The constants and path coefficients
introduced here are inputs; their simultaneous choice along a renormalization-group trajectory
will require the later estimates.

\section{Local gauge transformations and a symmetric choice}
\label{ch04:frames}

A choice of gauge on a patch assigns a frame to each site. Relative to
the original frames, this choice is specified by a family of matrices
$h_x\in G$: a contour variable $U_{xy}$ then has matrix $h_xU_{xy}h_y^{-1}$.
We choose $h_0=1$, so that $h_x$ expresses vectors at $x$ in the frame
at the patch centre. This lets us compare contour variables of the fine field before averaging. A deterministic axial gauge is convenient for bounds
but singles out coordinate directions. We combine a symmetric minimizing
choice near the flat field with a finite randomized choice elsewhere.
The use of gauge-covariant averages in a renormalization transformation belongs to
Ba\l{}aban's renormalization-group method~\cite{Balaban1987RG}; the precise
branch rules below are part of the present construction.

Let $\mathbb B_R=\{-R,\ldots,R\}^4$. For a permutation $o$ of the four
coordinate directions, denote by $T_U(x\to y;o)$ the contour variable along the
contour which changes coordinates in order $o$. For an initial point $z\in\mathbb B_R$,
put
\begin{equation}
 g_x^{z,o}=T_U(z\to x;o),\qquad
 h_x^{z,o}=(g_0^{z,o})^{-1}g_x^{z,o}.
 \label{ch04:axial-frame}
\end{equation}
These are \emph{local gauge transformations}, that is, $G$-valued functions on the sites of $\mathbb B_R$; they bring $U$ to an axial gauge with initial point $z$, and we call them \emph{axial}. The normalized one satisfies $h_0^{z,o}=1$. The unnormalized one gives
link coordinates
$U_{(x,\mu)}^{z,o}=g_x^{z,o}U(x,\mu)(g_{x+e_\mu}^{z,o})^{-1}$.
Their distances from the identity are unchanged if the normalized transformation is used instead,
since that change conjugates every transformed link by the same matrix.

A symmetric choice is possible if all relevant holonomies are close to the identity.
For $x,y\in\mathbb B_R$ and orders $o,o_0$, define the holonomy based at the
patch center by
\[
 \mathcal T_{xy;o}^{o_0}
   =g_x^{0,o_0}T_U(x\to y;o)(g_y^{0,o_0})^{-1}.
\]
Gauge transformations conjugate this holonomy, preserving its distance from the identity. Fix a small constant $\delta_{\rm fr}>0$.
If $\|\mathcal T_{xy;o}^{o_0}-1\|_G<\delta_{\rm fr}$ for every
$x,y\in\mathbb B_R$ and every $o,o_0\in S_4$, use the
normalized minimizer, with $h_0=1$, of
\begin{equation}
 \mathcal E_U(h)=\frac1{48|\mathbb B_R|^2}
   \sum_{x,y\in\mathbb B_R}\sum_{o\in S_4}
       \bigl\|h_xT_U(x\to y;o)-h_y\bigr\|_G^2.
 \label{ch04:frame-energy}
\end{equation}
The condition that all these distances be strictly below $\delta_{\rm fr}$, the \emph{loop condition}, includes every order $o_0$, so it does not privilege a preliminary
axial gauge.

If the loop condition fails, use axial gauge transformations. Fix $a_{\rm ax}\ge16$ and call a pair
$(z,o)$ eligible when
\begin{equation}
 \sum_{e\text{ internal to }\mathbb B_R}\|U_e^{z,o}-1\|_G
 \le a_{\rm ax}(2R+1)
      \sum_{\square\text{ internal to }\mathbb B_R}\|U_\square-1\|_G.
 \label{ch04:axial-eligible}
\end{equation}
Choose one pair $(z,o)\in\mathbb B_R\times S_4$ uniformly among all
eligible pairs. Use this same initial point and order throughout the patch:
the gauge at each $x\in\mathbb B_R$ is given by $h_x^{z,o}$ in
\eqref{ch04:axial-frame}. The coordinate paths from $z$ form a spanning
tree, and in this gauge every link of that tree is the identity. On the minimizing branch give all
pairs the same uniform probability weight, but let each pair specify the same minimizing
gauge transformation. The redundant values of the index $(z,o)$ allow both branches to be described by one
finite sum.

The following local assertion ensures that these rules are well-defined.
Its uniformity in the size of the patch is essential.

\begin{lemma}[Uniform choice of the local gauge transformation]
\label{ch04:frame-lemma}
There is a $\delta_{\rm fr}>0$, depending only on the matrix group, for
which the minimizer of \eqref{ch04:frame-energy} with $h_0=1$ is unique
whenever the loop condition holds. For every configuration and every fixed
coordinate order, at least one initial point satisfies \eqref{ch04:axial-eligible}
with $a_{\rm ax}=16$. These constants do not depend on $R$.
\end{lemma}
\begin{proof}
We first record the geometric input. The nearest-point projection $\pi$
onto the compact smooth submanifold $G$ is uniquely defined and smooth
on a uniform tubular neighborhood~\cite[p.~153, Lemma]{Foote1984}.
Its derivative at $1$ is the orthogonal projection $\Pi_{\gLie}$ onto the
Lie algebra. Hence on a sufficiently small convex ball about $1$,
\begin{equation}
 \|\pi(A)-\pi(B)-\Pi_{\gLie}(A-B)\|_G
       \le\tfrac14\|A-B\|_G.
 \label{ch04:projection-bound}
\end{equation}
This is the mean-value estimate for $D\pi-\Pi_{\gLie}$, not a new
projection theorem.

Fix $o_0\in S_4$ and use the axial gauge with initial point at the origin, with this
coordinate order. In this gauge the contour variable from $x$ to $y$ in order
$o$ is $\mathcal T_{xy;o}^{o_0}$. On the minimizing branch the loop condition
gives $\|\mathcal T_{xy;o}^{o_0}-1\|_G<\delta_{\rm fr}$.
Reversing a contour inverts its contour variable and permutes the coordinate orders.
A minimizer exists by compactness. Common left multiplication leaves the
energy invariant, so minimization with the normalization $h_0=1$ also gives
an unconstrained minimizer after allowing that common multiplier.
Comparison with the constant gauge transformation and the squared triangle inequality give
\[
 \frac{1}{|\mathbb B_R|^2}\sum_{x,y\in\mathbb B_R}\|h_x-h_y\|_G^2\le4\delta_{\rm fr}^2.
\]
Multiply all $h_x$ on the left by one group element
so that $1$ is nearest to their centroid $\bar h=|\mathbb B_R|^{-1}\sum_{x\in\mathbb B_R} h_x$.
Then $\|\bar h-1\|_G\le2\delta_{\rm fr}$. For each fixed $x$, minimization over the value at that point shows
that $h_x$ is nearest to
\[
 F_x(h)=\frac{1}{24|\mathbb B_R|}\sum_{y\in\mathbb B_R}\sum_{o\in S_4}h_y(\mathcal T_{xy;o}^{o_0})^{-1}.
\]
Indeed, when varying $h_x$ to $v$, retaining the old $h_x$ in the
diagonal term $y=x$ adds a nonnegative multiple of $\|v-h_x\|_G^2$,
which vanishes at $v=h_x$.
Since $\|F_x(h)-\bar h\|_G\le\delta_{\rm fr}$, all local means $F_x(h)$ lie
within $3\delta_{\rm fr}$ of $1$. For small enough $\delta_{\rm fr}$ they lie in the projection ball,
and $h_x=\pi(F_x(h))$.

Multiply two minimizers $h,k$ on the left in this way, so that $1$ is nearest to each centroid, and set
$u=\max_x\|h_x-k_x\|_G$. The differences of their local means satisfy
\[
 \|F_x(h)-F_x(k)\|_G\le u,\qquad
 \|F_x(h)-F_x(k)-(\bar h-\bar k)\|_G\le \sqrt{n_G}\,\delta_{\rm fr} u.
\]
For the second inequality, apply $\|AB\|_G\le\|A\|_G\|B\|_{\rm op}$
to each summand and use $\|B\|_{\rm op}\le\sqrt{n_G}\|B\|_G$.
The differences $\bar h-1$ and $\bar k-1$ are normal to $G$ at $1$, so
$\Pi_{\gLie}(\bar h-\bar k)=0$. At an index attaining the maximum,
\eqref{ch04:projection-bound} gives $u\le(\sqrt{n_G}\,\delta_{\rm fr}+1/4)u$.
Choosing $\sqrt{n_G}\,\delta_{\rm fr}\le1/4$
forces $u=0$. Thus minimizers differ only by a common left multiplier,
removed by $h_0=1$. Only averages occur in the contraction, which explains
why its constants do not grow with the patch.

For the axial gauge transformations, apply the usual plaquette telescoping
argument to a rectangular slice $X$ with $r$ varying coordinates. Write
$\Sigma_{\rm link}(X,U)$ and $\Sigma_{\rm plaq}(X,U)$ for the sums of
$\|U_e-1\|_G$ over its internal positive links and
$\|U_\square-1\|_G$ over its internal plaquettes, respectively. Choose a
hyperplane $X_0$ of the slice perpendicular to the last coordinate whose internal plaquette sum
is no greater than the average over such hyperplanes. Use the axial construction on
that hyperplane, then extend it along the last coordinate, setting those links
to one. Every remaining link differs from its translate, along the last coordinate, in the
chosen hyperplane by a product of intervening plaquettes. By the triangle
inequality and unitary invariance of $\|\cdot\|_G$, with $n=2R+1$ we obtain
\[
 \Sigma_{\rm link}(X,U^{\rm ax})
 \le n\bigl(\Sigma_{\rm link}(X_0,U^{\rm ax})+
                            \Sigma_{\rm plaq}(X,U)\bigr).
\]
Induction with the bound $4rn\Sigma_{\rm plaq}(X,U)$ closes because
$n\Sigma_{\rm plaq}(X_0,U)\le\Sigma_{\rm plaq}(X,U)$.
The construction follows the prescribed coordinate order; the chosen
hyperplanes of the successive slices, one in each dimension, specify an initial point. At $r=4$ the bound is \eqref{ch04:axial-eligible}
with constant $16$.
\end{proof}

Only smoothness of the nearest-point projection was used here. Analytic
estimates for the critical configuration will need more; they are not consequences of
this lemma alone.

\section{The gauge-covariant averaging operation and the exact kernel}
\label{ch04:kernel-section}

Translate the preceding construction of local gauge transformations to every patch
$Ly+\mathbb B_R$ centered at the fine site $Ly$ corresponding to the coarse site $y$, the \emph{patch of the averaging operation} at $y$.
For $z\in\mathbb B_R$, let $h_{y,z}$ be the gauge matrix at $Ly+z$,
normalized by $h_{y,0}=1$. For a coarse edge $e=(y,\mu)$ define the \emph{gauge-transformed parallel transporter}, for short the \emph{transporter},
\begin{equation}
 Y_{e,z}=h_{y,z}
       \left(\prod_{r=0}^{L-1}U(Ly+z+re_\mu,\mu)\right)
                     h_{y+e_\mu,z}^{-1}.
 \label{ch04:connector}
\end{equation}
Both ends are transformed by the local gauge transformations of the coarse endpoints, as shown
for two offsets in Figure~\ref{ch04:fig-block-paths}. In
particular, under a fine gauge transformation $g$, this transporter becomes
$g(Ly)Y_{e,z}g(Ly+Le_\mu)^{-1}$.

Let the \emph{averaging weight} $w^{\rm av}$ be the $j$-fold convolution of the uniform probability
on $\mathbb B_{(L-1)/2}$. It is nonnegative, has sum one, is invariant under
$H_4$, and is supported in $\mathbb B_{j(L-1)/2}\subset\mathbb B_R$.
Put $\bar Y_e=\sum_{z\in\mathbb B_R} w_z^{\rm av}Y_{e,z}$. On the first of its two branches, the value of the gauge-covariant \emph{averaging operation} is obtained by projecting the weighted
average of the transporters to the group; it then minimizes
\begin{equation}
 \mathcal P_e(u)=\tfrac12\sum_z w_z^{\rm av}\|Y_{e,z}-u\|_G^2,\qquad u\in G.
 \label{ch04:predictor-cost}
\end{equation}
Expanding about the weighted mean gives
\[
 2\mathcal P_e(u)
 =\sum_z w_z^{\rm av}\|Y_{e,z}-\bar Y_e\|_G^2
                         +\|\bar Y_e-u\|_G^2.
\]
Thus the minimizers are precisely the nearest group points to $\bar Y_e$.
If $\min_{u\in G}\mathcal P_e(u)<\varepsilon_{\rm pred}$, the identity
above gives $\min_{u\in G}\|\bar Y_e-u\|_G
<\sqrt{2\varepsilon_{\rm pred}}$. Choose $\varepsilon_{\rm pred}$ small
enough that $\bar Y_e$ then lies in the tubular neighborhood used in
Lemma~\ref{ch04:frame-lemma}. The value of the averaging operation is then unique, independently
of the number of transporters or the weights.

On this branch set $Q_e^{\rm pred}$ to that minimizer and assign uniform
probability weight to the redundant values $z\in\mathbb B_R$ of the index of the averaging operation. If the condition on the minimum fails,
let the index take the value $z$, and set $Q_e^{\rm pred}=Y_{e,z}$, with probability
$w_z^{\rm av}$. Equality in either the loop condition or the condition on
$\min_{u\in G}\mathcal P_e(u)$ selects the corresponding axial gauge transformation
or transporter. The averaging operation $U\mapsto Q_e^{\rm pred}$ is an analogue of Ba\l{}aban's averaging operation $\bar U$~\cite{Balaban1987RG}, not the same map: on the minimizing branch it is the nearest group point to a weighted mean, where Ba\l{}aban takes the exponential of a mean of logarithms.

\begin{figure}[htbp]
\centering
\begin{tikzpicture}[
    x=.85cm,y=.85cm,font=\small,
    transport/.style={draw=ymInk,line width=1.1pt,
      decoration={markings,mark=at position .64 with {\arrow{Stealth[length=2.6mm]}}},
      postaction={decorate}},
    frame/.style={draw=ymAccent,line width=.9pt,densely dashed,
      -{Stealth[length=2.2mm]},shorten <=2pt,shorten >=2pt}]
  \foreach \x in {0,...,8}
    \draw[ymGrid,thin] (\x,-1.3) -- (\x,2.3);
  \foreach \y in {-1,...,2}
    \draw[ymGrid,thin] (-.35,\y) -- (8.35,\y);

  % Two translated fine paths, each with seven bonds in the displayed slice.
  \draw[transport] (1,2) -- (8,2);
  \draw[transport] (0,-1) -- (7,-1);
  \foreach \r in {0,...,7}{
    \fill[ymInk] (1+\r,2) circle (1.5pt);
    \fill[ymInk] (\r,-1) circle (1.5pt);
  }
  \node[above=2pt] at (1,2) {$Ly+z$};
  \node[above=2pt] at (8,2) {$L(y+e_\mu)+z$};
  \node[above=2pt] at (4.5,2) {$L$ fine links};

  % Curves denote frame factors even when the frame is defined by minimization.
  \draw[frame] (0,0) to[bend left=16]
    node[pos=.50,left=3pt,fill=white,inner sep=1pt] {$h_{y,z}$} (1,2);
  \draw[frame] (8,2) to[bend left=16]
    node[pos=.50,right=3pt,fill=white,inner sep=1pt] {$h_{y+e_\mu,z}^{-1}$} (7,0);
  \draw[frame] (0,0) to[bend right=20] (0,-1);
  \draw[frame] (7,-1) to[bend right=20] (7,0);
  \foreach \x in {0,7}
    \fill (\x,0) ++(-2pt,-2pt) rectangle ++(4pt,4pt);
  \node[left=4pt,fill=white,inner sep=1pt] at (0,0) {$Ly$};
  \node[right=4pt,fill=white,inner sep=1pt] at (7,0) {$L(y+e_\mu)$};
\end{tikzpicture}
\caption[Fine paths and endpoint frames in the block average]{Two connectors in the matrix average $\bar Y_e$ for $e=(y,\mu)$,
shown in a coordinate plane. Each solid path traverses $L$ fine links.
The dashed arrows are frame factors, which can depend on the whole
patch. They give the framed transport $Y_{e,z}$ the same coarse
endpoints for every offset $z$.}
\label{ch04:fig-block-paths}
\end{figure}


Collect all choices made for the local gauge transformations and for the averaging operation in a finite \emph{auxiliary index} $\alpha$ and
denote their joint probability, the probability weight of $\alpha$, by $\lambda(U,\alpha)$. The choices for the local gauge transformations are
made first, followed by conditionally independent choices for the averaging operation. Thus
\begin{equation}
 \lambda(U,\alpha)\ge0,\qquad \sum_\alpha\lambda(U,\alpha)=1.
 \label{ch04:tag-law}
\end{equation}
Lemma~\ref{ch04:frame-lemma} supplies an eligible pair whenever it is needed.

Fix the \emph{penalty parameter} $b>0$ of the block kernel.\martin{Apparently $b$ will be fixed throughout and could just be
taken as $b=1$. It controls the soft conditioning used here.} The renormalization transformations for gauge fields in~\cite{Balaban1984Propagators,Balaban1985Averaging,Balaban1987RG} have a delta function in place of the densities $q_{b\beta}$ used below, so that there the coarse field equals the average exactly. For $\gamma\ge0$, let
\[
 q_\gamma(u)=Z_\gamma^{-1}e^{-\gamma\|u-1\|_G^2/2},\qquad
 Z_\gamma=\int_G e^{-\gamma\|v-1\|_G^2/2}\,dh(v).
\]
This is a probability density on $G$ with respect to Haar measure.

For a fixed fine field $U$ and choice $\alpha$, the value of the averaging operation
$Q_e=Q_e^{\rm pred}(U,\alpha)$ represents transport across the coarse
edge $e$. We impose agreement of $V_e$ with $Q_e$ by a soft constraint.
By unitary invariance its energy is $\tfrac12 b\beta\|V_e-Q_e\|_G^2$.
This penalizes the discrepancy in a gauge-invariant way, since $V_e$
and $Q_e$ transform by the same matrices at their endpoints. At large
$b\beta$ the conditional law thus concentrates near the value of the averaging operation.
Given $U$ and $\alpha$, we sample the coarse links independently with
these normalized densities. Averaging their product over $\alpha$ gives
the conditional density of $V$ given $U$; integrating against the fine
density gives its coarse marginal.

\begin{definition}[Exact block map]
\label{ch04:block-map}
The superscript $(L)$ denotes the block ratio: on a coarse torus
$\Lambda_s$ the fine torus is $\Lambda_{Ls}$. The period $s$, smoothing
order $j$ and penalty parameter $b$ are understood. For $\beta\ge0$, define the
block kernel of the step and its action on fine densities by
\begin{align}
 k^{(L)}_\beta(U,V)
   &=\sum_\alpha\lambda(U,\alpha)
        \prod_{e\text{ coarse}}
           q_{b\beta}\bigl(V_eQ_e^{\rm pred}(U,\alpha)^{-1}\bigr),
 \label{ch04:kernel}\\
 (\mathcal B^{(L)}_\beta f)(V)
   &=\int k^{(L)}_\beta(U,V)f(U)\,dU,
 \label{ch04:block-operator}
\end{align}
where $dU$ and $dV$ denote product probability Haar measure on their
respective link spaces.
\end{definition}

For fixed $\beta$ this is a bounded Borel kernel. Indeed, the loop
inequalities of Section~\ref{ch04:frames} and the axial eligibility
inequality \eqref{ch04:axial-eligible} compare continuous functions of
the fine links. Once the local gauge transformations are chosen, the condition for using the
minimizing branch of the averaging operation on edge $e$ is
$\min_{u\in G}\mathcal P_e(u)<\varepsilon_{\rm pred}$. This minimum is
continuous in the transporters $(Y_{e,z})_z$ by compactness of $G$.
On each minimizing branch the unique minimizer depends continuously on
its data. Thus the choices and their probabilities in the finite sum
\eqref{ch04:kernel} are Borel functions of $U$. Boundedness follows
from boundedness of $q_{b\beta}$ and \eqref{ch04:tag-law}.

\begin{proposition}[Normalization and covariance]
\label{ch04:kernel-properties}
The kernel is nonnegative and normalized in the coarse variable:
\begin{equation}
 \int k^{(L)}_\beta(U,V)\,dV=1.
 \label{ch04:kernel-normalization}
\end{equation}
It is covariant under fine gauge transformations and the induced coarse
gauge transformations $y\mapsto g(Ly)$. It is also covariant under signed
coordinate permutations and corresponding fine/coarse translations.
For every integrable fine density $f$,
\begin{equation}
 \int\mathcal B^{(L)}_\beta f\,dV=\int f\,dU.
 \label{ch04:exact-integral}
\end{equation}
\end{proposition}
\begin{proof}
For normalization, integrate each coarse-link factor using right invariance
of Haar measure, then sum \eqref{ch04:tag-law}. Fubini gives
\eqref{ch04:exact-integral}, with absolute integrability following from
the same computation applied to $|f|$.
For covariance, gauge transformations conjugate the holonomies of the loop conditions and
transform the endpoints of \eqref{ch04:connector}; signed lattice
isometries permute initial points, orders and offsets and invert reversed links.
The objective functions and eligibility conditions are unchanged. Uniqueness
on the minimizing branches, and reindexing of the finite alternatives,
then give the stated covariance of the kernel.
\end{proof}

In particular, if $U$ has normalized density $f/\int f$, the coarse
field has density $\mathcal B^{(L)}_\beta f/\int f$. The equality of
integrals \eqref{ch04:exact-integral} holds even when the coarse action
has terms beyond the Wilson plaquette action.

\section{A linear bound on coarse plaquette distances}
\label{ch04:predictor-geometry}

We bound the distance of a coarse plaquette holonomy from the identity
in terms of the fine plaquette holonomies in a surrounding box and the
discrepancies between the four coarse boundary links and the values of
the averaging operation on them. Fix a site $y$ of the coarse torus
$\Lambda_s$. Choose a lift $\tilde y\in\ZZ^4$ and form the fine lattice box
\[
 M_y=L\tilde y+\mathbb B_{(j+7)L}.
\]
For a fine plaquette $q$ on $\Lambda_{Ls}$, let $m_y(q)$ be the number
of its lifts with all four vertices in $M_y$. Changing the lift of $y$
translates the whole box by a fine period and leaves this count unchanged.
If $s>2(j+7)$, the box can contain at most one lift of each plaquette. Write
\[
 \mathcal P_y(U)=\sum_{q}m_y(q)\|U_q-1\|_G,
\]
where the sum ranges over all fine plaquettes on the torus, with one
orientation per square.

\begin{lemma}[Coarse plaquette estimate]
\label{ch04:predictor-comparison}
Fix the smoothing order $j$ and the constants in the choices of the local
gauge transformations and of the averaging operation in
Sections~\ref{ch04:frames} and~\ref{ch04:kernel-section}.
Suppose $s\ge100(j+8)$. Let $p$ be a coarse plaquette whose four
vertices lie within periodic sup distance $2$ of $y$ in coarse lattice
units. In particular, $y$ may be any vertex of $p$. For every fine field
$U$, coarse field $V$, and choice $\alpha$ with $\lambda(U,\alpha)>0$,
write $Q_e(U)=Q_e^{\rm pred}(U,\alpha)$. Then
\begin{equation}
 \|V_p-1\|_G\le\sum_{e\in\partial p}\|V_e-Q_e(U)\|_G
                         +C_jL^{-2}\mathcal P_y(U),
 \qquad \sup_q\sum_{y\in\Lambda_s}m_y(q)\le(2j+15)^4 .
 \label{ch04:macro-comparison}
\end{equation}
The constant $C_j$ is independent of $L$ and $s$.
\end{lemma}

\begin{lemma}[Mean distance from a constant gauge transformation]
\label{ch04:minimizing-frame-bound}
Let $U$ be a link field on a finite lattice box $X$, and let $h$ minimize
the objective \eqref{ch04:frame-energy} on $X$. Take uniform averages
over the indicated sites $x,z\in X$ and coordinate orders $o\in S_4$.
There exists $c\in G$ such that
\begin{equation}
 \operatorname{av}_x\|h_x-c\|_G
 \le C_G\operatorname{av}_{x,z,o}\|T_U(x\to z;o)-1\|_G ,
 \label{ch04:frame-first-power}
\end{equation}
where $C_G$ depends only on $G$.
\end{lemma}
\begin{proof}
Put $d_x=\operatorname{av}_{z,o}\|T_U(x\to z;o)-1\|_G$ and
$\bar d=\operatorname{av}_x d_x$. The constant gauge transformation has energy at most $\bar d$, since each path distance
is at most two. Minimality and the squared triangle inequality give
$\operatorname{av}_{x,z}\|h_x-h_z\|_G^2\le8\bar d$.
Let $m=\operatorname{av}_z h_z$ and
$m_x=\operatorname{av}_{z,o}h_zT_U(x\to z;o)^{-1}$, so
$\|m_x-m\|_G\le d_x$. As in the proof of Lemma~\ref{ch04:frame-lemma},
varying $h_x$ shows that it is a nearest group element to $m_x$.
Indeed, reversing the contour interchanges the two off-diagonal terms
involving $x$. Freezing the second endpoint of the diagonal term adds
a nonnegative squared distance that vanishes at $h_x$.

Choose $c$ nearest to $m$. Since each $h_x$ lies in $G$, Jensen's
inequality gives $\|m-c\|_G^2\le\operatorname{av}_x\|m-h_x\|_G^2
\le8\bar d$. Choose $\varepsilon>0$ so that the nearest-point
projection is Lipschitz on the ball of radius $2\varepsilon$ about $c$,
uniformly in $c\in G$, as in Lemma~\ref{ch04:frame-lemma}.
First suppose $8\bar d<\varepsilon^2$. At sites with $d_x<\varepsilon$
both $m$ and $m_x$ lie in this ball, so $\|h_x-c\|_G\le C_Gd_x$.
At the remaining sites use $\|h_x-c\|_G\le2\le2d_x/\varepsilon$.
Averaging gives \eqref{ch04:frame-first-power}. For
$8\bar d\ge\varepsilon^2$ the bound follows from $\|h_x-c\|_G\le2$.
The same proof applies on any finite
set for a uniformly averaged family of paths paired under reversal
with identity contour variables on the diagonal.
\end{proof}

\begin{proof}[Proof of Lemma~\ref{ch04:predictor-comparison}]
Let $X$ be a fine lattice box with $\ell$ sites in each direction and
$N=\ell^4$. For a link field $U$ on $X$, define the sum over its internal
positive links by
\[
 E_X(U)=\sum_{\mu=1}^4\ \sum_{\substack{x\in X\\x+e_\mu\in X}}
                       \|U(x,\mu)-1\|_G .
\]
For each coordinate order $o$ we have
\[
 \sum_{x,z\in X}\|T_U(x\to z;o)-1\|_G\le N\ell E_X(U).
\]
Indeed, choose $x,z$ independently and uniformly in $X$. On the
segment in direction $\mu$, the other three coordinates are uniform on
the $\ell^3$ transverse sites. Each is taken from $x$ or $z$ according
to the order $o$. The triangle inequality bounds the segment distance by
the sum of link distances along its entire coordinate line. Its expectation
is therefore at most $\ell^{-3}$ times the sum over all positive
$\mu$-links in $X$. Sum over $\mu$ and multiply by $N^2$.

Let $H$ be any gauge transformation on $X$. If $f$ is the minimizing local gauge transformation for $U$,
then $h_x=f_xH_x^{-1}$ minimizes the objective \eqref{ch04:frame-energy} for $H\cdot U$,
since substituting $h$ gives the same energy term by term. A common
left multiplication can restore the normalization at the centre.
Applying Lemma~\ref{ch04:minimizing-frame-bound} and the path bound gives
\[
 \sum_{x\in X}\|f_x-cH_x\|_G\le C_G\ell E_X(H\cdot U)
\]
for some $c\in G$. For an eligible axial gauge transformation $f$, put $g_x=f_xH_x^{-1}$.
The pure gauge $g\cdot1$ has contour variable $g_xg_z^{-1}$ from $x$ to $z$,
and the linkwise triangle inequality gives
$E_X(g\cdot1)\le E_X(f\cdot U)+E_X(H\cdot U)$. Applying the path
bound to $g\cdot1$ and averaging over $z$ yields some $c=g_z$ with
\[
 \sum_{x\in X}\|f_x-cH_x\|_G
 \le \ell\bigl(E_X(H\cdot U)+E_X(f\cdot U)\bigr).
\]
By the eligibility condition \eqref{ch04:axial-eligible},
$E_X(f\cdot U)\le a_{\rm ax}\ell\sum_{q\subset X}\|U_q-1\|_G$,
where the sum runs over plaquettes with all four vertices in $X$.

Now choose the axial gauge transformation $H$ of Lemma~\ref{ch04:frame-lemma} on $M_y$.
Since this box has side $2(j+7)L+1$, the axial estimate gives
$E_{M_y}(H\cdot U)\le C_j L\mathcal P_y(U)$. Lift the four vertices
$z$ of $p$ to $\tilde y+\mathbb B_2$. Their patches of the averaging operation
$Lz+\mathbb B_R$, with $R=(j+3)L$, and the straight paths between them
all lie in $M_y$. Each patch has side $2R+1=O_j(L)$. Applying the
preceding bounds on that patch, for either choice of local gauge transformation, gives $c_z\in G$
such that
\begin{equation}
 \sum_{i\in\mathbb B_R}\|h_{z,i}-c_zH_{Lz+i}\|_G
                                  \le C_j L^2\mathcal P_y(U).
 \label{ch04:patch-frame-sum}
\end{equation}
The averaging weights satisfy $w_i^{\rm av}\le L^{-4}$. Given the
last $j-1$ uniform points, at most one of the $L^4$ possible first
points yields a prescribed sum. Thus each endpoint local gauge transformation in the
transporter \eqref{ch04:connector} contributes at most
$C_jL^{-2}\mathcal P_y(U)$ to the weighted discrepancy. For its middle
contour, at each of the $L$ steps translation of the offset $i$ is injective.
The unweighted sum of the path distances in gauge $H$ is therefore at
most $L E_{M_y}(H\cdot U)$. Combining these bounds gives, for each link
of $\partial p$ written in its positive orientation $e=(z,\mu)$,
\[
 \sum_i w_i^{\rm av}\|Y_{e,i}-c_zc_{z+e_\mu}^{-1}\|_G
                                      \le C_jL^{-2}\mathcal P_y(U).
\]

For either branch of the averaging operation and any $u_*\in G$ we have
\[
 \|Q_e(U)-u_*\|_G\le2\max\{1,\varepsilon_{\rm pred}^{-1}\}
                      \sum_i w_i^{\rm av}\|Y_{e,i}-u_*\|_G .
\]
On the minimizing branch $Q_e(U)$ is a nearest group point to the
weighted mean of the $Y_{e,i}$, so the triangle inequality gives the
factor two. On the other branch the minimum of
\eqref{ch04:predictor-cost} is at least $\varepsilon_{\rm pred}$.
Every group distance is at most two, so
$\varepsilon_{\rm pred}\le\mathcal P_e(u_*)\le
\sum_i w_i^{\rm av}\|Y_{e,i}-u_*\|_G$ and $\|Q_e(U)-u_*\|_G\le2$.
Take $u_*=c_zc_{z+e_\mu}^{-1}$ on each link.
These four transports multiply to the identity around $p$, with
inverses on reversed links. Replacing one factor at a time, first by
the corresponding $Q_e(U)$ and then by $V_e$, proves the first
inequality in \eqref{ch04:macro-comparison}.

For the overlap bound, fix one lift of $q$ with base vertex $a\in\ZZ^4$.
Summing $m_y(q)$ over coarse sites is equivalent to fixing this lift
and counting all centres $Lz$, $z\in\ZZ^4$, whose boxes contain it.
Indeed, a lift translated by a fine period is in $M_y$ exactly when the
fixed lift is in the oppositely translated box. Each coordinate must
satisfy $|Lz_\mu-a_\mu|\le(j+7)L$, leaving at most $2(j+7)+1=2j+15$
integer choices for $z_\mu$. The four coordinates give the stated
bound $(2j+15)^4=O(j^4)$.
\end{proof}

\section{The regularity condition and the representation of the density}
\label{ch04:density-coordinates}

The threshold used to distinguish small from large fields must remain fixed
while the inverse coupling $\beta$ is varied for tuning. To express that distinction,
let $\bar\beta>1$ and use the reference values
\begin{align}
 \widehat\beta_m&=\bar\beta+\kappa_G\tau_m,&
 M_m^{\rm cut}&=\exp\bigl((\log\log\widehat\beta_m)^2\bigr),\\
 t_m&=\widehat\beta_m^{-1/2}M_m^{\rm cut},&
 B_m&=\lceil(\log\widehat\beta_m)^2\rceil.
 \label{ch04:thresholds}
\end{align}
Here $\tau_m=\log D_m$ is the logarithmic scale from
Definition~\ref{ch01:scales}. We use these formulas only where $0<t_m\le1$.
The inverse coupling $\beta$ ranges in
$\mathfrak B_m=[\widehat\beta_m/2,2\widehat\beta_m]$.
The reference value and threshold are not recomputed when that coupling
changes. Their eventual compatibility with the schedule of Chapter~\ref{ch10:trajectory} is a later
choice.

At the present scale write $\Lambda=\Lambda_{R_pD_m}$, and let
$\operatorname{Bad}_m(U)$, the \emph{large-field set} of $U$, be the sites incident to at least one plaquette
with $\|U_\square-1\|_G>t_m$. The configuration $U$ satisfies the \emph{regularity condition} on a set $X$, which is then called a small-field region, if
$X\cap\operatorname{Bad}_m(U)=\varnothing$; write the indicator of this condition as $h_X(U)$.
Thus every plaquette meeting $X$ must have $\|U_\square-1\|_G\le t_m$. The term regularity condition is that of~\cite{Balaban1985Spaces}, where the spaces of regular configurations carry bounds on the plaquette variables, depending on the scale, on a nested sequence of domains, and a second condition; here the condition is a single inequality at one scale, and the large-field set is neither enlarged nor summed over. This definition is
gauge invariant and involves a closed condition on the sites outside the large-field set.

We keep two kinds of local corrections: the localized terms, whose sum is the \emph{regular part of the action}, and the large-field activities.

\begin{definition}[Localized terms and large-field activities]
\label{ch04:activities}
A \emph{localized term} of the regular part of the action consists of a \emph{localization domain} $X_i$, a base point $o_i\in X_i$,
and a real Borel gauge-invariant function $E_i^{\rm reg}$ local on $X_i$.
It vanishes at the flat configuration and has the required number of
ordered left derivatives on configurations satisfying the regularity condition on $X_i$.

A \emph{large-field activity} has an \emph{exclusion support} $Y_i$, a nonempty \emph{prescribed large-field
set} $I_i^{\rm bad}\subset Y_i$, and an additional Borel gauge-invariant
condition, an event $\mathcal E_i^{\rm extra}$. Its real Borel value
$K_i^{\rm lf}$ depends on links internal to $Y_i$ and on the indicators of
large-field sites in $Y_i$, and vanishes outside
\begin{equation}
 \{U:\operatorname{Bad}_m(U)\cap Y_i=I_i^{\rm bad}\}
                         \cap\mathcal E_i^{\rm extra}.
 \label{ch04:eligibility}
\end{equation}
The additional condition reads the same links and large-field indicators.
\end{definition}
The names localized term, localization domain and regular part of the action are Ba\l{}aban's~\cite{Balaban1987RG,Balaban1988Expansions}. His localized terms are analytic, one for each connected localization domain, and bounded term by term; those used here are real Borel functions with finitely many derivatives, have a base point and an arbitrary localization domain, and are summed in a weighted $\ell^1$ norm. The word regular has two technical uses, as in his papers: the regular part of the action, and the regularity condition. Ba\l{}aban uses the word activity for the activities of a polymer expansion~\cite{Balaban1988Cluster}; its use for the large-field terms is particular to this book. Neither a localized term nor a large-field activity is assumed nonnegative. Both may be indexed by a
countable set; their summability will be part of the norm conditions.

There are two different meanings of support in this definition. The set
$Y_i$ determines exclusion and coverage: the $Y_i$ of each subfamily in the sum below will be pairwise disjoint.
But deciding whether a site belongs to $\operatorname{Bad}_m(U)$ reads
plaquettes meeting that site, whose links can extend beyond $Y_i$. The
\emph{dependence set} includes those links as well as the internal
links and any links read by other conditions. Two disjoint exclusion supports
can have overlapping dependence sets. A connected expansion based
on dependence must use the latter, whereas the sum over disjoint covering families below uses
the former. Disjoint exclusion supports do not imply probabilistic
independence.

\section{The prescribed leading action}
\label{ch04:leading-action}

The running coupling is a component of the representation relative to a prescribed leading
action. Its construction will come from Gaussian integration and local
matching; here we specify the form of that component without claiming
the estimates needed to select it.

At scale $m$, fix constants $c_{\rm dir},c_{\rm lin},R^{\rm rel}>0$ with
$16c_{\rm lin}<c_{\rm dir}$. Define
\begin{align}
 \operatorname{clip}_t(r)&=
 \begin{cases}
 c_{\rm dir}r^2/2,&r\le t,\\
 c_{\rm lin}tr,&r>t,
 \end{cases}
 \label{ch04:clip}\\
 \chi^{\rm rel}_{m,x}(U)&=
 \mathbf1\left\{\|U_{y;\mu\nu}-1\|_G
       \le t_m e^{\operatorname{dist}(x,y)/R^{\rm rel}}
              \text{ for every }y,\mu<\nu\right\}.
 \label{ch04:relaxed-test}
\end{align}
The distance is periodic and measured in lattice units. We call the condition in \eqref{ch04:relaxed-test} the \emph{relaxed regularity condition} at $x$: the regularity condition on the whole lattice, with $t_m$ replaced by the threshold $t_m e^{\operatorname{dist}(x,y)/R^{\rm rel}}$ growing exponentially with the distance from $x$. Although written
as a condition on every plaquette, \eqref{ch04:relaxed-test} reads only a finite
neighborhood: $\|U_{y;\mu\nu}-1\|_G\le2$ makes all conditions with
$t_m e^{\operatorname{dist}(x,y)/R^{\rm rel}}\ge2$ automatic. Any claim of
locality must include the plaquettes of the remaining conditions in its dependence
set.

For each orientation $\mathfrak p=(\mu,\nu)$, fix a summable linear
combination of transported curvatures,
\[
 \mathcal L^{\rm path}_{\mathfrak p}(U,x)
      =\sum_i r_{\mathfrak p i}\mathcal F_{\mathfrak p i}(U,x).
\]
Each $\mathcal F_{\mathfrak p i}$ is a plaquette curvature transported to
$x$, as in \eqref{ch02:transported-curvature}. Also specify local real
invariant multilinear expressions $A_i^{(q)}(U,x)$, $q=2,3,4$, in such
curvatures, each multiplied by the indicator of the regularity condition on its declared
support. The support of the expression contains the contours, plaquettes and base point; this
indicator may read additional adjacent plaquettes. Assume the coefficient sums
converge absolutely at the fixed finite volume. Set
\begin{align}
 P^{\rm raw}_{N,m}(\beta,U)
 &=\beta\sum_{x,\mathfrak p}
    \left[\operatorname{clip}_{t_m}(\|U_{x;\mathfrak p}-1\|_G)
       +\tfrac12\chi^{\rm rel}_{m,x}(U)
                    |\mathcal L^{\rm path}_{\mathfrak p}(U,x)|^2\right]
 \notag\\
 &\quad+\beta\sum_{x,i}(A_i^{(3)}+A_i^{(4)})(U,x)
                     +\sum_{x,i}A_i^{(2)}(U,x),\\
 P^{\rm path}_{N,m}(\beta,U)
 &=\frac1{|H_4|}\sum_{h\in H_4}P^{\rm raw}_{N,m}(\beta,h^*U).
 \label{ch04:path-action}
\end{align}
The coefficients are held fixed as $\beta$ varies. Thus this action is
affine in $\beta$. Its ultimate normalization requires the quadratic
coefficient in the $\beta$-dependent part to be the curvature square,
and no such coefficient in the $\beta$-independent quadratic term $\sum_{x,i}A_i^{(2)}(U,x)$. Those are
matching conditions, not consequences of the displayed form.

The function $\operatorname{clip}_t$ of \eqref{ch04:clip}, quadratic below $t$ and linear above, is not a smooth cutoff.
At $r=t$ it drops from $c_{\rm dir}t^2/2$ to $c_{\rm lin}t^2$ immediately
to the right. The terms of the representation of the effective density are therefore Borel functions,
not globally differentiable functions of the links. Derivative estimates
refer to the local expressions on the domains where the regularity condition holds, with the event indicators
held fixed. In particular, a derivative with respect to the inverse coupling $\beta$
does not differentiate the threshold or its indicator.

\section{Representations by disjoint covering families, and their norms}
\label{ch04:density-norms}

Recall that $Y_i$ is an exclusion support and $I_i^{\rm bad}$ its prescribed large-field
set, as in Definition~\ref{ch04:activities}. Let
$\mathfrak F_m(U)$ be the finite subsets $\mathfrak a$ of large-field
activity indices such that the $Y_i$, $i\in\mathfrak a$, are pairwise
disjoint and
$\bigcup_{i\in\mathfrak a}I_i^{\rm bad}=\operatorname{Bad}_m(U)$.
We call them the \emph{disjoint covering families} of the large-field sites (in combinatorics, exact covers). They do not select a
preferred partition into connected components.

\emph{Representation data} $\rho$ consist of $\beta$, a real constant
$v_{\rm sc}$ per site, the localized terms and the large-field
activities, with the fixed action \eqref{ch04:path-action}. Define
\begin{align}
 V_\rho(U)&=P^{\rm path}_{N,m}(\beta,U)
                +\sum_i h_{X_i}(U)E_i^{\rm reg}(U)-v_{\rm sc}|\Lambda|,
 \notag\\
 \operatorname{dens}(\rho)(U)
   &=e^{-V_\rho(U)}
          \sum_{\mathfrak a\in\mathfrak F_m(U)}
                         \prod_{i\in\mathfrak a}K_i^{\rm lf}(U).
 \label{ch04:effective-density}
\end{align}
The empty product is one. Because each $I_i^{\rm bad}$ is nonempty, if
there are no large-field sites the only disjoint covering family is empty and the sum
is one. Otherwise a single support may cover two regions, or two disjoint
supports may cover them separately; both alternatives belong in the sum.
Representation data satisfying the definitions alone need not give a positive density.
Positivity follows when a later construction proves that its density equals
an exact block integral of a positive input.

Fix an integer derivative order $q\ge2$, a nonnegative spatial exponent
$\varpi$, a support exponent $\sigma>0$ (the exponent of the support weight), and the threshold $t_m$.
For a localized term with base point $o_i$, use the seminorm
\begin{equation}
 \|E_i^{\rm reg}\|_m
 =\sup_{0\le k\le q}\ \sup_{U,A_1,\ldots,A_k}
        |D_{A_1}^L\cdots D_{A_k}^L E_i^{\rm reg}(U)|,
 \label{ch04:local-seminorm}
\end{equation}
where $U$ ranges over configurations that satisfy the regularity condition on $X_i$.
Each $A_r$ is a Lie algebra valued link field satisfying, on each link $e$ with initial site $e^-$,
$|(A_r)_e|_{\gLie}\le t_m(1+\operatorname{dist}(e^-,o_i))^\varpi$.
The left derivatives are those of Section~\ref{ch02:bounds}. The required
differentiability includes existence and uniform bounds at the boundary of the
domain defined by the regularity condition, interpreted through the local expression $E_i^{\rm reg}$ itself;
no derivative of $h_{X_i}$ is included.

For a set $X\subset\Lambda$, let the \emph{connected covering number} $\operatorname{cov}_{B_m}(X)$ be the
minimum number of radius-$B_m$ periodic balls whose union is connected
in the nearest-neighbor graph and covers $X$. Set $\operatorname{cov}_{B_m}(\varnothing)=0$. The covering union may
contain additional sites. Its exponential weight penalizes extended
supports without requiring a support itself to be connected; the covering number plays the role of the linear size of a localization domain in~\cite{Balaban1987RG}.
Choose majorants $e_i\ge\|E_i^{\rm reg}\|_m$ and
$k_i\ge\sup_U|K_i^{\rm lf}(U)|$, and define
\begin{equation}
 \begin{split}
 \mathcal E_m(\rho)&=\sup_{x\in\Lambda}
     \sum_{i:o_i=x}e^{\sigma\operatorname{cov}_{B_m}(X_i)}e_i,\\
 \mathcal K_m(\rho)&=\sup_{x\in\Lambda}
     \sum_{i:x\in Y_i}e^{\sigma\operatorname{cov}_{B_m}(Y_i)}k_i.
 \end{split}
 \label{ch04:weighted-sums}
\end{equation}
The first is a sum over the terms with base point $x$, an \emph{anchored} sum: each localized term is counted at one
site. The second counts a large-field activity at every site of its
exclusion support. Replacing an anchored sum by a sum over the terms whose localization domain contains $x$, a \emph{containing-site} sum like the second,
requires a bound for the number of possible base points; it is not a harmless
notational change. Likewise, replacing $Y_i$ by its dependence set
changes the weight and requires a geometric estimate.

Finite weighted sums ensure that \eqref{ch04:effective-density} is
well-defined at fixed volume. Indeed
\[
 \sum_i\|E_i^{\rm reg}\|_m\le|\Lambda|\mathcal E_m,\qquad
 \sum_i\sup|K_i^{\rm lf}|\le|\Lambda|\mathcal K_m,
\]
and the absolute sum over all finite subfamilies is at most
$\prod_i(1+k_i)\le\exp(|\Lambda|\mathcal K_m)$.
Restricting to disjoint covering families only lowers this bound. The estimate
is a finite-volume convergence statement, not a uniform bound on the
entire partition function.

\section{Conditions for a one-step induction}
\label{ch04:admissibility}

The preceding representation data describe possible densities; their smallness
conditions will describe the inputs of a one-step theorem. There are
additional normalization conditions. On a nonwrapping localization domain\mgubi{Issue 10: ``nonwrapping'' is used in several chapters and never defined. In the sources: $10\,\mathrm{dist}(x,o)<n$ for all $x$ in the support, $n$ the period; Chapter 17 uses the factor $10$.}, choose
centered coordinate lifts $\xi_{o_i}(y)$. For every pair of affine link
fields of the form
\[
 A_{(y,\mu)}=u_\mu+\sum_\nu\xi_{o_i}(y)^\nu v_{\mu\nu},
 \qquad u_\mu,v_{\mu\nu}\in\gLie,
\]
require the Hessian at zero of $A\mapsto E_i^{\rm reg}(\exp A)$ to vanish
on that pair. Together with $E_i^{\rm reg}(1)=0$ and the vanishing linear
term, this excludes the low-order local contributions assigned to the
leading action. The first derivative vanishes for a gauge-invariant localized
term by invariance under constant gauge transformations: each
linear coefficient would otherwise be an invariant linear functional on
the simple Lie algebra. The meaning of the affine Hessian condition is
fixed in the chosen local coordinates; no coordinate lift is imposed on
a wrapping localization domain by this definition.

One may prescribe the intended size bounds
\begin{equation}
 \mathcal E_m\le\widehat\beta_m^{-11/10},\qquad
 \mathcal K_m\le\exp\bigl(-(M_m^{\rm cut})^{1/5}\bigr),
 \label{ch04:activity-bounds}
\end{equation}
with the reference quantities of \eqref{ch04:thresholds}. In an induction,
$q,\varpi,\sigma$ may depend on the \emph{stage} $J_m$ in force at scale $m$ (a stage specifies finitely many derivative orders, support exponents, path widths and coefficient bounds; see Section~\ref{ch10:choices}); the stage index
is different from the scale index $m$. For now they remain specified input
parameters. Neither \eqref{ch04:activity-bounds} nor the density identity
proves that the block map preserves the class.

Parameter continuity is also stated for the representation data. A curve
$u\mapsto\rho_u$ has fixed localization domains, exclusion supports, base points, prescribed large-field sets,
additional conditions and event indicators. Its coupling and its constant $v_{\rm sc}$ are continuous. At each
$u_0$, differences of the derivatives of the localized terms and of the values of the large-field activities have
majorants $e_i(u),k_i(u)$ satisfying
\[
 \sum_i e^{\sigma\operatorname{cov}_{B_m}(X_i)}e_i(u)
 +\sum_i e^{\sigma\operatorname{cov}_{B_m}(Y_i)}k_i(u)
                  \longrightarrow0\quad(u\to u_0).
\]
Both sums here are over all indices, at a fixed finite volume. This
continuity condition is strong enough to control the density and its
integral; it differentiates no moving large-field boundary.

The analytic task remaining is now precise. Starting with representation data at the
fine scale, one must construct coarse representation data whose density is exactly
$\mathcal B^{(L_m)}_\beta\operatorname{dens}(\rho)$, whose coupling has the
claimed adjustment, and whose localized terms and large-field activities obey the coarse conditions.
A class defined as the outputs reachable from a Wilson starting density
can only be introduced after that construction. It must not be confused
with the much larger set of representation data satisfying the displayed
norm bounds.
